\section{INTRODUÇÃO}
\label{sec:introduction}

Estabelecido em 2007 com o Decreto n\(^\circ\) 6.022/2007 \cite{decreto_6022},
o Sistema Público de Escrituração Digital, abreviado como SPED, é um sistema que
busca modernizar a interação entre contribuintes e a administração tributária
pública encarregada de regulamentar, arrecadar, fiscalizar e auditar a prestação
de contas e o pagamento dos tributos, administração essa denominada ao longo
desse documento como \textit{fisco}. O SPED executa essa modernização através da
informatização das obrigações acessórias, bem como a utilização da certificação
digital para garantir a validade jurídica dos documentos eletrônicos
\cite{sped_modernizacao}. Um dos módulos desse sistema é a Escrituração Fiscal
Digital de ICMS --- Imposto sobre Operações relativas à Circulação de
Mercadorias e sobre Prestações de Serviços de Transporte Interestadual e
Intermunicipal e de Comunicação --- e IPI --- Imposto sobre Produtos
Industrializados ---, abreviada como EFD ICMS IPI, que opera por meio de um
arquivo digital, cuja emissão é obrigatória para contribuintes desses impostos,
contendo registros de apuração de impostos, documentos fiscais e outras
informações de interesse da Secretaria da Receita Federal, dos fiscos estaduais
e do Distrito Federal \cite{efd_icms_ipi}.

O documento EFD ICMS IPI deve ser escriturado pelo contribuinte,
validado se utilizando do programa especificado pelos fiscos das unidades
federadas e pelo fisco federal, e submetido a programa validador, disponilizado
em sítio digital do SPED, para validação de conteúdo, assinatura digital e
transmissão \cite{guia_pratico_efd}. A responsabilidade pela governança
e custódia desses dados recai diretamente sobre a Administração Tributária.
Sob a égide do Decreto n\(^\circ\) 6.022/2007 \cite{decreto_6022}, o sistema
é explicitamente definido em seu Artigo 2\(^\circ\) como ``\textit{[...]
instrumento que unifica as atividades de recepção, validação, armazenamento e
autenticação de livros e documentos que integram a escrituração contábil e
fiscal dos empresários e das pessoas jurídicas, inclusive imunes ou isentas,
mediante fluxo único, computadorizado, de informações.}''. Portanto, a
manutenção íntegra e perpétua desse acervo digital configura um dever legal
imperativo do Estado, essencial para o exercício pleno da fiscalização soberana.

Além disso, no âmbito da Lei 5.172/1966 \cite{lei-5172}, em seu
Artigo 150\(^\circ\), Parágrafo 4\(^\circ\), é explicitado que
``\textit{Se a lei não fixar prazo a homologação, será êle de cinco anos, a
contar da ocorrência do fato gerador; expirado êsse prazo sem que a Fazenda
Pública se tenha pronunciado, considera-se homologado o lançamento e
definitivamente extinto o crédito, salvo se comprovada a ocorrência de dolo,
fraude ou simulação}'', sendo também reforçado pelo artigo 173\(^\circ\) desse
mesmo decreto, que afirma ``\textit{O direito de a Fazenda Pública constituir o
crédito tributário extingue-se após 5 (cinco) anos, contados: I --- do primeiro
dia do exercício seguinte àquele em que o lançamento poderia ter sido efetuado;
II --- da data em que se tornar definitiva a decisão que houver anulado, por
vício formal, o lançamento anteriormente efetuado}'', de modo que além do
armazenamento de tais documentos, é de interesse do fisco realizar auditorias
sobre esses documentos em tempo hábil, para evitar fraudes fiscais e
homologação involuntárias nos termos desses decretos, o que pode gerar
prejuízos bilionários aos cofres públicos.

Contudo, a efetivação dessas auditorias esbarra em diversos desafios de
engenharia de software e processamento de dados. Os arquivos da EFD ICMS IPI
são semiestruturados, apresentando uma limitação física de blocos, mas deixando
de fora a hierarquia entre os registros. Essa hierarquia, bem como outras
regras de formatação, formam o que é denomidado de \textit{leiaute}, o qual
sofre atualizações frequentes por parte do fisco \cite{nota_tecnica_efd}.
A ingestão desse volume massivo de dados semiestruturados para sistemas
gerenciadores de bancos de dados (SGBD) relacionais --- utilizando abordagens
tradicionais de código fixo --- resulta em alto custo de manutenção e
dificuldades no mapeamento de entidades, como limitações na manipulação e de
cruzamento nesses grandes volumes de dados.

Com este cenário estabelecido, o presente trabalho se propõe a desenvolver um
\textit{pipeline} de ingestão dinâmico e orientado a metadados para o
processamento de arquivos EFD ICMS IPI. A solução busca abstrair as regras de
extração do código-fonte, por meio de arquivos de configuração em formato JSON
(Notação de Objetos do JavaScript, tradução livre do inglês
\textit{JavaScript Object Notation}), permitindo que a infraestrutura do fisco
processe e armazene os dados fiscais de forma confiável e resiliente às
constantes mudanças de leiaute.

\subsection{Objetivo geral}
\label{sec:introduction_objetivo_geral}
Desenvolver uma arquitetura de ingestão e processamento de dados orientada a
metadados para arquivos EFD ICMS IPI, visando permitir o armazenamento
estruturado e facilitar o cruzamento de dados fiscais.

\subsection{Objetivos específicos}
\label{sec:introduction_objetivos_especificos}
Para que o objetivo geral seja alcançado, foram traçados os seguintes objetivos específicos:
\begin{itemize}
    \item Determinar, dentre as diversas regras especificadas na documentação
    oficial \cite{guia_pratico_efd, nota_tecnica_efd}, quais são relevantes
    para o nosso objetivo;
    \item Mapear e estruturar as regras de negócio relevantes dos leiautes do
    SPED Fiscal utilizando arquivos de metadados estruturado em JSON;
    \item Implementar uma ferramenta que seja capaz de ler, extrair e atribuir
    a estrutura do leiaute aos blocos de texto da EFD ICMS IPI de forma
    dinâmica, \textit{i.e.}, sem dependência de código fixo;
    \item Validação da integridade dos dados.
\end{itemize}

\subsection{Estrutura do relatório técnico}
\label{sec:introduction_estrutura}

Para uma melhor compreensão e fluidez da leitura, as próximas páginas se
encontram divididas em um total de 5 seções, além das referências
bibliográficas. A Seção 2 aborda trabalhos semelhantes ao presente. A Seção 3,
por sua vez, detalha com mais riqueza o contexto do problema, usando isso como
base para apresentar a solução elaborada. A Seção 4 expõe os resultados obtidos
e faz uma análise de cada um deles, cobrindo tópicos específicos que foram
considerados relevantes. A Seção 5 reúne as considerações finais, as
contribuições da pesquisa e as sugestões para trabalhos futuros. Por fim, a
Seção 6 apresenta a declaração formal sobre o uso de ferramentas de Inteligência
Artificial ao longo da elaboração deste documento.